\section{The SDP minion}
\label{sec:minion}


As previously mentioned, polymorphisms are a powerful tool for understanding the computational complexity of PCSPs. However, beyond some of the simplest classes of PCSPs, individually classifying the complexity based on specific polymorphisms can be unwieldy. Instead, one often looks to higher-level structure between classes of polymorphisms, which is captured by the notion of \emph{minions} (also called clonoids) \cite{BBKO21}.

In this section, we describe the structure of the ``basic SDP minion,'' which gives a precise algebraic condition for when the basic SDP decides a Promise CSP. We note that a similar minion was concurrently and independently discovered by Ciardo-Zivny \cite{cz22-minion}. 

{We prove our minion characterization for ``ordinary'' promise constraint satisfaction problems over any domain. In particular, we do not consider instances with folding or the setting of constants.}

\subsection{Minion preliminaries}






We now give the definition of a minion. {This is most easily stated in the language of category theory, however we shall also give ``concrete'' interpretations of any category-theoretic statements we make. We let $\Set$ denote the category of sets, and $\NeFinSet$ denote the category of \emph{non-empty} finite sets.}

\begin{definition}
A minion $\mathcal M$ is {a functor from $\NeFinSet$ to $\Set$. More concretely, for every non-empty finite set $A$, we let $\cM^A$ denote the image of $A$ with respect to the functor. For every pair of non-empty finite sets $A$ and $B$ and function $\pi : A \to B$, we let the \emph{minor map} $\cM^{\pi} : \cM^{A} \to \cM^{B}$ denote the image of $\pi$ with respect to the functor. Further, for any maps $\pi : A \to B$ and $\tau : B \to C$, we have that $\cM^{\tau} \circ \cM^{\pi} = \cM^{\tau \circ \pi}$. We also assert for the identity map $\id_A : A \to A$, we have that $\cM^{\id_A} : \cM^A \to \cM^A$ is also the identity map.}
\end{definition}

To streamline the minor map notation $\cM^{\pi}$, for all $M \in \cM^A$, we let $\pi \oslash M := \cM^{\pi}(M) \in \cM^B$.

{
\begin{remark}
Some works only use \emph{minion} (or \emph{clonoid}) to refer to sets of the form $\Pol(\Gamma)$ for some (possibly) promise template $\Gamma$, and refer to all other minions as an \emph{abstract minions}. In this work, all such objects are minions.

In addition, $\cM$ is typically only defined as a functor from the subcategory of $\NeFinSet$ of sets of the form $[k]$ for some $k \in \N$. Our definition is equivalent as every non-empty finite set $A$ has a bijection with $[|A|]$. However, we consider this richer definition of a minion so that we can discuss how the minion acts on sets without an explicit enumeration. 
\end{remark}
}

\subsubsection{{Examples}}

{We now walk through a few examples.}
%
\paragraph{Polymorphisms.} {Given a promise template $\Gamma = \{(P_1, Q_1), \hdots, (P_l, Q_l)\}$ over the domain $(D_1, D_2)$, one can view $\Pol(\Gamma)$ as a minion. Formally, for every non-empty finite set $A$, we let $\Pol(\Gamma)^A$ denote the set of functions $f : D_1^A \to D_2$ (where $D_1^A$ is the set of functions from $A$ to $D_1$) such that for any $i \in [l]$ and $\alpha : A \to P_i$, we have that 
\begin{align}
    (f(\pi_1 \circ \alpha), \hdots, f(\pi_{k_i} \circ \alpha)) \in Q_i,\label{eq:pol-cond}
\end{align}
where $\pi_j : P_i \to D_1$ is the projection of $P_i \subseteq D^{k_i}$ onto the $j$th coordinate. Observe that if $A = \{1,\hdots, L\}$ then $\Pol(\Gamma)^A$ can be identified with the arity $L$ polymorphisms.

Likewise, for any map $\tau : A \to B$ for non-empty finite sets $A$ and $B$, we define $\Pol(\Gamma)^\tau$ such that for any $f \in \Pol(\Gamma)^A$  and $\beta : B \to D_1$, we have that
\begin{align}
    (\tau \oslash f)(\beta) = f(\beta \circ \tau).\label{eq:pol-min}
\end{align}
We It's clear that $\id_{A} \oslash f = f$. Furthermore, for any pair of maps $\tau : A \to B$, $\tau' : B \to C$, and $\gamma : C \to D_1$, we have that
\[
  (\tau' \oslash (\tau \oslash f))(\gamma) = (\tau \oslash f)(\gamma \circ \tau') = f((\gamma \circ \tau') \circ \tau)) = f(\gamma \circ (\tau' \circ \tau)) = ((\tau' \circ \tau) \oslash f)(\gamma),  
\]
so associativity also holds.
}

{
\paragraph{Trivial minion.} The simplest example of a minion $\cM_{\operatorname{triv}}$ maps every non-empty finite set to the same set $\{e\}$. All functions $\tau : A \to B$ map to the same trivial map from $e$ to istelf.
}

{
\paragraph{Identity minion.} {The \emph{identity minion} $\cM_{\id}$ is the identity functor from $\NeFinSet$ to itself (when viewed as a subcategory of $\Set$). In other words, for every non-empty finite set $A$ we have that $\cM_{\id}^A = A$, and for every map $\tau : A \to B$, we have that $\cM_{\id}^\tau = \tau$.}
}

{
\paragraph{Convex combinations.} For every non-empty finite set $A$, let $\Delta_A$ be the set of probability distributions over $A$. That is, $\Delta_A = \{f : A \to [0,1] : \sum_{a\in A} f(a) = 1\}.$ We let $\cQ_{\conv}$ denote the minion which maps $A$ to $\Delta_A$~\cite{BBKO21}. Given a map $\tau : A \to B$ and a probability distribution $f \in \Delta_A$, we let $\tau \oslash f \in \Delta_B$ be defined by 
\[
    (\tau \oslash f)(b) = \sum_{a \in \tau^{-1}(b)} f(a),
\]
where the empty sum is equal to $0$. The main motivation for studying $\cQ_{\conv}$ is that it is closely related to the use of basic LP \cite{BBKO21}.
}

\subsubsection{{Minion homomorphisms}}

{
A key approach of \cite{BBKO21} is that to better understand $\PCSP(\Gamma)$, we can ask how $\Pol(\Gamma)$ relates to other (abstract) minions. The most natural way to do that is via a \emph{minion homomorphism}, which in category-theoretic language is a \emph{natural transformation} between two minions.}

\begin{definition}
A \emph{minion homomorphism} $\psi: \mathcal M \to \mathcal N$ between two minions {is a natural transformation between the respective functors. More concretely, for every non-empty finite set $A$, we have a map $\psi_A : \cM^A \to \cN^A$ such that for every map $\tau : A \to B$ and $M \in \cM^A$, we have that}
\[
    {\tau \oslash \psi_A(M) = \psi_B(\tau \oslash M).}
\]
In other words, the following diagram commutes:
\begin{center}
\begin{tikzcd}
    \cM^A \arrow[r,"\psi_A"] \arrow[d,"\cM^{\tau}"] &\cN^A \arrow[d,"\cN^{\tau}"]\\
    \cM^B \arrow[r,"\psi_B"] &\cN^B
\end{tikzcd}
\end{center}
\end{definition}

{To assert the existence of a minion homomorphism from $\mathcal M$ and $\cN$, we write $\cM \to \cN$. A key result of \cite{BBKO21} is that minion homomorphisms govern the computational complexity of PCSPs.}

{\begin{theorem}[\cite{BBKO21}]
Let $\Gamma_1, \Gamma_2$ be promise templates. If $\Pol(\Gamma_1) \to \Pol(\Gamma_2)$ then decision version of $\PCSP(\Gamma_2)$ is polynomial-time reducible to $\PCSP(\Gamma_1)$.
\end{theorem}}

{Another result of \cite{BBKO21} is that for any promise template $\Gamma$, $\cQ_{\conv} \to \Pol(\Gamma)$ if and only if the basic LP solves $\PCSP(\Gamma)$. Our goal in the remainder of this section is to prove an analogous result for the basic SDP.} 

\subsubsection{The free structure}

{Let $\Gamma$ be a promise template.} One of the most important tools for understanding minion {homomorphisms of} the form $\mathcal M \to \Pol({\Gamma})$ {is} the \emph{free structure} \cite{BBKO21}. {Since in this paper we parameterize our promise constraint satisfaction problems using promise templates rather than a pair of relational structures, we restate the definition of a free structure in our language, which we call a \emph{free template}.}  {Given a relation $R \subseteq D^k$, we define the \emph{free relation} $\F_{\cM}(R)$ to be a predicate of arity $k$ over $\cM^D$ with the following property, we have that $(M_1, \hdots, M_k) \in (\cM^D)^k$ is an element of $R$ if and only if there exists $N \in \cM^{R}$ such that for all $i \in [k]$, we have that 
\[
    \pi_i \oslash N = M_i,
\]
where $\pi_i : R \to D$ is the projection of $R$ onto the $i$th coordinate.

Given a promise template $\Gamma = \{(P_1, Q_1), \hdots, (P_l, Q_l)\}$, we define the \emph{free template} 
\[
\F_{\cM}(\Gamma) := \{(\F_{\cM}(P_1), Q_1), \hdots, (\F_{\cM}(P_l), Q_l)\}.
\]
Recall that for a promise template $\Gamma$ over domain $(D_1, D_2)$, there exists a map $h : D_1 \to D_2$ such for every $i \in [\ell]$ and $\bx \in P$, we have that $h(\bx) \in Q$. It is not obvious (nor necessarily true) that for $\F_{\cM}(\Gamma)$ an analogous map $h^{\cM} : \cM^{D_1} \to D_2$ necessarily exists. In fact, constructing $h^{\cM}$ asserts the existence of a minion homomorphism. 
}
{This is t}he fundamental property of the free structure{.}

{
\begin{lemma}[Lemma~4.4 of \cite{BBKO21}]\label{lem:free-structure}
$\mathcal M \to \Pol(\Gamma)$ if and only if $\F_{\cM}(\Gamma)$ is a promise template.
\end{lemma}
}

{In the language of category theory, the free structure can be thought of as an adjoint functor. See the discussion preceding Remark 4.5 in \cite{BBKO21}.} 


\subsection{SDP {m}inion {d}efinition}

{
We now define the minion corresponding to the behavior of the basic SDP. We call this minion $\cM_{\SDP}.$ Recall that for $\cQ_{\conv}$, for any non-empty finite set $A$, we have that $\cQ_{\conv}^A = \Delta_A$, where $\Delta_A$ is the set of probability distributions over $A$. We define $\cM_{\SDP}^A$ in an analogous manner which captures orthogonal configurations of vectors index by $A$.}

{S}ince the SDP minion is a universal object, we need to be able to represent vectors of arbitrary large dimensions. We achieve this using infinite-dimensional vectors that are eventually zero. Similar techniques have been used in other minion constructions \cite{ciardo2022clap}.

Let $\R^{\omega}$ be infinite sequences of real numbers {indexed by $\N$} which are eventually $0${.} {We note} that $\R^{\omega}$ is an inner product space {with the bilinear form}
\[
    {\langle \bx, \by \rangle := \sum_{i=1}^{\infty} x_i y_i,}
\]
{which is well-defined as only finitely many terms in the summation are nonzero.}

{We now formally define $\cM_{\SDP}$. For each non-empty finite set $A$, we define $\cM_{\SDP}^{A}$ to be the set of functions $\bw : A \to \R^{\omega}$ with the following properties.}

\begin{enumerate}
\item For all {distinct $a, a' \in A$, $\langle \bw(a), \bw(a')\rangle = 0$.}
\item $\sum_{{a \in A}} \|{\bw(a)}\|^2_2 = 1$.
\end{enumerate}
Observe that the {first condition implies that the} second condition is equivalent to $\|\sum_{{a \in A}} {\bw(a)}\|_2 = 1$.

The minors of $\cM_{\SDP}$ are {defined analogously to $\cQ_{\conv}$. Let $A$ and $B$ be non-empty finite sets and consider $\bw \in \cM_{\SDP}^A$ and $\tau : A \to B$. We define $\tau \oslash \bw : B \to \R^{\omega}$ to be for all $b \in B$,}
\[
    {(\tau \oslash \bw)(b) := \sum_{a \in \tau^{-1}(b)} \bw(a),}
\]
{where the empty sum equals $\mathbf{0} \in \R^{\omega}$. We now finish the remaining details which prove that $\cM_{\SDP}$ is a minion.}

{\begin{lemma}
$\cM_{\SDP}$ is a minion.
\end{lemma}}
\begin{proof}
First, for each {$\bw \in \cM_{\SDP}^{A}$} and {$\tau : A \to B$} we verify that ${\bw' := \tau \oslash \bw \in \cM_{\SDP}^{B}}$. First, fix {distinct $b, b' \in B$}. We have that
\begin{equation*}
    { \langle \bw'(b) , \bw'(b') \rangle }= { \left\langle \Bigl(\sum_{a \in \tau^{-1}(b)} \bw(a)\Bigr), \Bigl(\sum_{a' \in \tau^{-1}(b')} \bw(a')\Bigr)\right\rangle}
    = \sum_{{\substack{a \in \tau^{-1}(b)\\a' \in \tau^{-1}(b')}}} { \langle \bw(a) , \bw(a')\rangle}
    = 0{,}
\end{equation*}
{since $b$ and $b'$ are distinct and $\tau$ is a function.} Further,
\begin{align*}
{ \left\langle \Bigl( \sum_{b \in B} \bw'(b)\Bigr), \Bigl(\sum_{b \in B} \bw'(b)\Bigr) \right\rangle} &= { \left\langle \Bigl(\sum_{a\in A} \bw(a)\Bigr), \Bigl(\sum_{a\in A} \bw(a)\Bigr) \right\rangle}= 1.
\end{align*}
Thus, ${\bw' \in \cM_{\SDP}^{B}}.$ It is trivial to observe that {$\cM^{\id_A}$ is the identity map.}

The only remaining condition to check is that the minors commute. Consider {$\pi : A \to B$} and {$\tau : B \to C$. Consider $\bu \in \cM_{\SDP}^A.$}  We seek to verify that $\tau \oslash (\pi \oslash \bu) = (\tau \circ \pi) \oslash \bu$. For all {$c \in C$}, we have that
{
\begin{align*}
(\tau \oslash (\pi \oslash \bu))(c) = \sum_{b \in \tau^{-1}(c)} (\pi \oslash \bu)(b) = \sum_{b \in \tau^{-1}(c)} \sum_{a \in \pi^{-1}(b)} \bu(a) = \sum_{a \in (\tau \circ \pi)^{-1}(c)} \bu(a) = ((\tau \circ \pi) \oslash \bu)(a),
\end{align*}
as desired.
}
\end{proof}

The goal of the rest of this section is to prove the following theorem.

\begin{theorem}\label{thm:basic-SDP}
{For any promise template $\Gamma$, t}he basic SDP decides $\PCSP({\Gamma})$ if and only if $\cM_{\SDP} \to \Pol({\Gamma})${.}
\end{theorem}

\subsection{On the basic SDP}

{As we are considering the decision version of the basic SDP, we assume that $\eps_j = 0$ for each clause index $j \in [m]$. For completeness, we state the simplified SDP below. Note that we omit an objective since we are concerned only with feasibility. Any undefined notation is the same as in Section~\ref{sec:prelims}. Given a clause index $j \in [m]$, we let $\cF_j$ denote the set of assignments $f : S_j \to D$ for which $f(C_j) \in P^{(j)}$.}
{
\begin{align*}
\text{(a1)} && \lambda_j(f) &\geq 0  & \forall j\in [m], ~ f \in \cF_j \\ 
\text{(a2)} && \sum_{f \in \cF_j}\lambda_j(f)&=1  & \forall j \in [m] \\
\text{(a3)} && \langle \bv_0, \bv_0\rangle &= 1\\
\text{(a4)} && { \langle \textbf{v}_{i,a},   \textbf{v}_0 \rangle } &= \sum_{\substack{f \in \cF_j\\f(x_i) = a}}\lambda_j(f)&  \forall j\in [m],~x_i \in C_j,~ a \in D \\ 
\text{(a5)} && { \langle \textbf{v}_{i,a} , \textbf{v}_{i',a'} \rangle } &= \sum_{\substack{f \in \cF_j\\f(x_i) = a, f(x_{i'}) = a'}}\lambda_j(f)  & \forall j\in [m],~x_i,x_{i'} \in C_j,~ a,a' \in D 
\end{align*}
}

{We call this formulation of the basic SDP the ``traditional'' basic SDP, in contrast to the ``alternative'' basic SDP we soon present.}

\begin{remark}
Although we specify here that the vectors can have an arbitrarily large dimension, it is known that an SDP with $n$ vector or scalar variables is feasible if and only if it is feasible in $n$ dimensions. Thus, nothing is lost from Section~\ref{sec:prelims} by making the vectors have an arbitrarily large (but finite) dimension.
\end{remark}

{\begin{remark}
Some formulations of the basic SDP also impose that for every pair of vectors appearing in the basic SDP has a nonnegative pairwise dot product. Raghavendra's theorem does not use this in their basic SDP formulation, so we do not include it in our formulation.
\end{remark}}

\subsubsection{An alternative basic SDP}

{In order to prove Theorem~\ref{thm:basic-SDP}, w}e now present a modified basic SDP which is more convenient to work with. {Instead of having an LP variable $\lambda_j(f)$ for each potential clause assignment $f : S_j \to D$, we have a vector $\bw_{j,f} \in \R^{\omega}.$ For each $i \in [n]$, let $\bv_i : D \to \R^{\omega}$ denote the function $\bv_i(a) = \bv_{i,a}$. Likewise, for every $j \in [m]$, let $\bw_j : \cF_j \to R^{\omega}$ be the function $\bw_j(f) = \bw_{j,f}$. With this notation, we can succinctly impose our SDP conditions as follows.}

{
\begin{align*}
\text{(b1)} &&  \bv_i &\in \cM_{\SDP}^{D} & \forall i \in [n] \\ 
\text{(b2)} && \bw_j &\in \cM_{\SDP}^{\cF_j}   & \forall j \in [m] \\ 
\text{(b3)} && \bv_{i,a} &= \sum_{\substack{f \in \cF_j\\f(x_i) = a}} \bw_{j,f} & \forall j \in [m], i \in S_j, a \in D,
\end{align*}
}
where we use the fact that membership in $\cM_{\SDP}^S$ can be expressed in $O(|S|^2)$ SDP constraints.

\begin{remark}
{Note that the special ``truth'' vector $\bv_0$ is omitted from this formulation. As we shall soon see, this does not change the power of the SDP relaxation, but it does make the analysis more straightforward.}
\end{remark}

{
\begin{remark}\label{rem:b3}
Note that (b3) can also be rewritten as $\bv_i = \tau_{i,j} \oslash \bw_j$, where $\tau_{i,j} : \cF_j \to D$ is the map $\tau_{i,j}(f) = f(x_i)$. As such, $\cM_{\SDP}$ can be swapped with \emph{any} minion $\cM$ in conditions (b1--b3). However, such a relaxation is not algorithmically tractable in general.
\end{remark}
}

\subsubsection{Equivalence with traditional basic SDP}

We now prove that these two formulations of the basic SDP are equivalent. 
That is, a solution to either SDP can be transferred to the other. 

{\begin{lemma}\label{lem:SDP-equiv}
For any promise template $\Gamma$ and any instance of $\PCSP(\Gamma)$, the traditional basic SDP for the instance is feasible if and only if the alternative basic SDP for the instance is feasible.
\end{lemma}}

{Before we begin, we note that we often apply linear maps to the vectors in order to transform vectors from one SDP solution to another, extending the ideas in Section~\ref{subsec:gram-ortho}. We let $\R^{\omega \times \omega}$ denote the set of matrices such that each row and column has finite support. These finite support conditions are such that for any $M \in \R^{\omega \times \omega}$ and $v \in \R^{\omega}$ we have that $Mv \in \R^{\omega}$ and that $\R^{\omega \times \omega}$ is closed under the transpose operator.}

{The identity matrix $I_{\omega} \in \R^{\omega \times \omega}$ has $(I_{\omega})_{i,j} = 1$ if $i = j$ and $0$ otherwise. We say that $Q \in \R^{\omega \times \omega}$ is \emph{orthogonal} if $Q Q^{T} = Q^T Q = I_{\omega}.$ We now extend Proposition~\ref{prop:gram-ortho} to our setting.}

{\begin{proposition}\label{prop:ortho}
Let $n$ be a positive integer and $\bu_1, \hdots, \bu_n, \bv_1, \hdots, \bv_n \in \R^{\omega}$ be vectors. We have that $\langle \bu_i, \bu_j \rangle = \langle \bv_i, \bv_j\rangle$ for all $i,j \in [n]$ if and only if there exists an orthogonal matrix $Q \in \R^{\omega \times \omega}$ such that $\bv_i = Q \bu_i$ for all $i \in [n]$. 
\end{proposition}
\begin{proof}
Note that for any $\bu, \bu' \in \R^{\omega}$ and $Q \in \R^{\omega \times \omega}$ orthogonal, we have that $\langle \bu, \bu'\rangle = \langle Q\bu, Q\bu' \rangle$. This proves the ``if'' direction. We now focus on proving the converse.

Assume that $\langle\bu_i, \bu_j\rangle = \langle \bv_i, \bv_j\rangle$ for all $i,j \in [n]$. Let $N$ be a positive integer such that $(\bu_i)_j = 0$ and $(\bv_i)_j = 0$ for all $i \in [n]$ and $j > N$. We can thus assume that $\bu_1, \hdots, \bu_n, \bv_1, \hdots, \bv_n \in \R^N$. By Proposition~\ref{prop:gram-ortho} there exists an orthogonal matrix $Q \in \R^{N \times N}$ such that $Q\bu_i = \bv_i$ for all $i \in [n]$. We can extend this $Q$ to an orthogonal $Q' \in \R^{\omega \times \omega}$ in the following straightforward manner.
\[
    Q'_{i,j} = \begin{cases}
    Q_{i,j} & i, j \le N\\
    1 & i = j > N\\
    0 & \text{otherwise}
    \end{cases}.\qedhere
\]
\end{proof}
}

\medskip\textbf{Alternative to traditional.} First, we show that a solution to the alternative basic SDP is a solution to the traditional basic SDP. For each ${i \in [n]}$, define ${\bv_{i,D} = \sum_{a \in D} \bv_{i,a}}$. By (b1), we know that each ${\bv_{i,D}}$ is a unit vector. Further by (b3), we can deduce that ${\bv_{i,D} = \bv_{i',D}}$ whenever ${x_i}$ and ${x_{i'}}$ appear in a common constraint ${C_j}$. Thus, ${\bv_{i,D} = \bv_{i',D}}$ whenever ${x_i}$ and ${x_{i'}}$ are in the same connected component of the {hypergraph with vertices $\{x_1, \hdots, x_n\}$ and hyperedges $\{C_1, \hdots, C_m\}$}. {Let $I \subseteq [n]$ and $J \subseteq [m]$ be the indices of the variables and constraints appearing in some connected component. For any orthogonal matrix $Q \in \R^{\omega \times \omega}$,} each of the constraints (b1), (b2), and (b3) are preserved when we {multiply each vector $\bv_{i,a}$ for $i \in I$ and $\bw_{j,f}$ for $j \in J$ on the left by $Q$.} Thus, by {multiplying by suitable orthogonal matrices (i.e., apply Proposition~\ref{prop:ortho} for $n=1$)}, we can achieve a solution to (b1-3) such that {each $\bv_{i,D}$} is the same unit vector for all ${i \in [n]}$, which we shall call ${\bv_0}$. {We call this an \emph{aligned} solution to the alternative basic SDP.}

{We now transform this aligned solution into a solution to the traditional basic SDP. Each $\bv_{x,a}$ in the traditional basic SDP is actually the same as it appeared in the aligned solution to the alternative basic SDP. We set $\bv_0$ to the newly-defined $\bv_0$. Finally, we set $\lambda_j(f) = \|\bv_{j,f}\|_2^2$ for all $f \in \cF_j$.}

To see that the traditional basic SDP is satisfied, note that condition (a1) follows by inspection, and condition (a2) follows from (b2). {Conditions (a3) and (a4) follow} by substituting ${\bv_{i,D}}$ to ${\bv_0}$ and applying (b3). Condition {(a5)} follows {by} combining (b2) and (b3).

\medskip\textbf{Traditional to alternative.} We now show that any solution to the traditional basic SDP is a solution to the alternative basic SDP. Assume we have a traditional basic SDP solution. Recall that ${\bv_{i,a}} \in \R^\omega$ {for each $i \in [n]$ and $a \in D$}. We shall use the same vectors ${\bv_{i,a}}$ in {the solution to the alternative basic SDP}. Note that (b1) then follows from {(a5)} and (a2). 

Since there are finitely many vectors, there exists some $N \in \N$ such that each ${\bv_{i,a}}$ {(and $\bv_0$)} has its support within the first $N$ coordinates. {For each $j \in [m]$ and $f \in \cF_j$, p}rovisionally assign ${\hat{\bw}_{j,f}}$ to be $\sqrt{\lambda_{{j}}({f})} \cdot e_i$, where {$e_i \in \R^{\omega}$ is the $i$th standard unit vector for some $i > N$} chosen uniquely for each constraint-assignment pair. As there are finitely many constraint-assignment pairs, all of these {new vectors} are in $\R^{\omega}$.

However, as written the ${\hat{\bw}_{j,f}}$'s are not compatible with the ${\bv_{i,a}}$'s. {To rectify this, for each $j \in [m], i \in S_j,$ and $a \in D$ we define}
\[{\hat{\bv}_{j,i,a}} := \sum_{\substack{{j \in \cF}\\{f(x_i)} = a}} {\hat{\bw}_{j,f}}\]
{Using (a5), we can readily verify that for any $j \in [m]$, $i,i' \in S_j$, and $a,a' \in D$, we have that $\langle \bv_{i,a}, \bv_{i,a'}\rangle = \langle \hat{\bv}_{j,i,a}, \hat{\bv}_{j,i',a'}\rangle$. Thus, by Proposition~\ref{prop:ortho}, for each $j \in [m]$, there exists an orthogonal matrix $Q_j \in \R^{\omega \times \omega}$ such that $Q_j \hat{\bv}_{j,i,a} = \bv_{i,a}$ for all $i \in S_j$ and $a \in D$.}

{For each $j \in [m]$ and $f \in \cF_j$, define $\bw_{j,f} := Q_j \hat{\bw}_{j,f}$. Observe that for each $j \in [m]$, $i \in S_j,$ and $a \in D$, we have that
\[
    \bv_{i,a} = Q_j \hat{\bv}_{j,i,a} = \sum_{\substack{{j \in \cF}\\{f(x_i)} = a}} Q_j {\hat{\bw}_{j,f}} = \sum_{\substack{{j \in \cF}\\{f(x_i)} = a}} \bw_{j,f}, 
\]}
as desired. Thus, conditions (b2) and (b3) hold. {Therefore,} the two SDP formulations are equivalent{, proving Lemma~\ref{lem:SDP-equiv}.}

\subsection{From minion homomorphism to SDP rounding algorithm}

Recall that we are trying to prove that {for any promise template $\Gamma$,} the basic SDP decides $\PCSP({\Gamma}) $ if and only if $\cM_{\SDP} \to \Pol({\Gamma})$.  We begin by showing the {easier} direction that the minion homomorphism implies that the basic SDP decides $\PCSP({\Gamma})$. {In other words, for any instance of $\PCSP(\Gamma)$, the existence of a vector solution to alternative basic SDP of the instance implies that there exists a weakly satisfying assignment to the instance.}

\begin{theorem}\label{thm:hom_to_alg}
If $\cM_{\SDP} \to \Pol({\Gamma})$, then the basic SDP decides $\PCSP({\Gamma})$.
\end{theorem}
{
\begin{proof}
Assume that $(D_1, D_2)$ is the domain of $\Gamma$ with homomorphism $h : D_1 \to D_2$. Consider an instance of $\PCSP(\Gamma)$ on variables $V = \{x_1, \hdots, x_n\}$ and clauses $\{C_1, \hdots, C_m\}$ such that for all $j \in [m]$, the clause $C_j$ corresponds to $(P^{(j)}, Q^{(j)}) \in \Gamma$ of arity $k_j$. Let $S_j$ be the set of variables appearing at least once in $C_j$. We also define $\cF_j$ as the set of assignments $f : S_j \to D_1$ for which $f(C_j) \in P^{(j)}$ and $\cG_j$ as the set of assignments $g : S_j \to D_2$ for which $g(C_j) \in Q^{(j)}$.

By Lemma~\ref{lem:SDP-equiv}, we may without loss of generality assume that there exist vectors $\{\bv_{i,a} : i \in [n], a \in D_1\}$ and $\{\bw_{j,f} : j \in [m], f \in S_j\}$ such that (b1), (b2), and (b3) hold, where we define $\bv_i(a) = \bv_{i,a}$ and $\bw_j(f) = \bw_{j,f}$.
If no such vectors exist, we know that our instance is not strongly satisfiable by the contrapositive of Proposition~\ref{prop:basic-SDP-completeness}.

Let $\psi : \cM_{\SDP} \to \Pol(\Gamma)$ be the guaranteed minion homomorphism. Let $\id_{D_1} : D_1 \to D_1$ be the identity map. We define our weak assignment $\sigma : V \to D_2$ as follows. For all $i \in [n]$, we set
\[
    \sigma(x_i) := \psi(\bv_i)(\id_{D_1}).
\]
Unpacking this, $\psi(\bv_i) \in \Pol(A, B)^{D_1}$, so $\psi(\bv_i)$ has type $\psi(\bv_i) : D_1^{D_1} \to D_2$. The ``canonical'' element of $D_1^{D_1}$ is the identity map, so we plug that into the function to get a value in $D_2$.

We now seek to show that for every $j \in [m]$, that $\sigma|_{S_j} \in \cG_j$. First, we need to eliminate any repetition of variables. To do this, observe that $(\cF_j, \cG_j)$ can be viewed as a promise relation of arity $S_j$ since it is obtained from $(P^{(j)}, Q^{(j)})$ by setting some tuples in the relation equal. As such, we have that $\Pol({\Gamma}) \subseteq \Pol(P^{(j)}, Q^{(j)}) \subseteq \Pol(\cF_j, \cG_j)$. Therefore, we can think of $\psi$ as a minion homomorphism from $\cM_{\SDP}$ to $\Pol(\cF_j, \cG_j)$. Thus, since $\bw_j \in \cM_{\SDP}^{\cF_j}$, we have that $\psi(\bw_j) \in \Pol(\cF_j, \cG_j)^{\cF_j}$. Thus, $\psi(\bw_j)$ has type $\cF_j^{D_1} \to D_2$.

For $x_i \in S_j$, let $\pi_i : \cF_j \to D_1$ be the evaluation map $\pi_i(f) = f(x_i)$. We check that for all $x_i \in S_j$ and $a \in D_1$, we have that
\begin{align}
(\pi_i \oslash \bw_{j})(a) = \sum_{\substack{f \in \cF_j\\f(x_i) = a}} \bw_{j,f} = \bv_{i}(a),\label{eq:pi-map}
\end{align}
by definition of the alternative basic SDP, so $\pi \oslash \bw_j = \bv_i$. Thus, for all $x_i \in S_j$ we have that
\begin{align*}
    \sigma(x_i) &= \psi(\bv_i)(\id_{D_1}) & \text{(definition of $\sigma$)}\\
    &= \psi(\pi_i \oslash \bw_j)(\id_{D_1})  & \text{(\ref{eq:pi-map})}\\
    &= (\pi_i \oslash \psi(\bw_j))(\id_{D_1}) & \text{($\psi$ is minion homomorphism)}\\
    &= \psi(\bw_j)(\id_{D_1} \circ \pi_i) & \text{(\ref{eq:pol-min})}\\
    &= \psi(\bw_j)(\pi_i \circ \id_{\cF_j}),
\end{align*}
where for the last line $\id_{D_1} \circ \pi_i = \pi_i = \pi_i \circ \id_{\cF_j}$.

Apply (\ref{eq:pol-cond}) with $A = \cF_i$ and $(P_i, Q_i) := (\cF_j, \cG_j)$. Since $\psi(\bw_j) \in \Pol(\cF_j, \cG_j)^{\cF_j}$ and $\id_{\cF_j} : \cF_j \to \cF_j$, we have that
\[
    (\sigma(x_i) : x_i \in S) = (\psi(\bw_j)(\pi_i \circ \id_{\cF_j}) : x_i \in S) \in \cG_j,
\]
as desired. Thus, $\sigma$ is indeed a satisfying assignment to the original instance. Therefore, the basic SDP decides $\PCSP(\Gamma)$.
\end{proof}
}

\subsection{From SDP rounding algorithm to minion homomorphism}
We now prove the converse. {To do this, we adapt} the proof techniques of {Theorem 4.11 of }\cite{ciardo2022clap}. 

\begin{theorem}
If the basic SDP decides $\PCSP({\Gamma})$, then $\cM_{\SDP} \to \Pol({\Gamma})$.
\end{theorem}

\begin{proof}
{Assume that $(D_1,D_2)$ is the domain of $\Gamma = \{(P_1,Q_1), \hdots, (P_l, Q_l)\}$ with a homomorphism $h : D_1 \to D_2$ such that $h(P_i) \subseteq Q_i$ for all $i \in [l]$. To simplify the proof, we make the following technical assumption similar to the reduction in Remark~\ref{rem:Ragh}. If a relation $(P,Q) \in \Gamma$ of arity $k$, then every possible relation which can be created by repeating variables is also in $\Gamma$. For example, if $P, Q$ have arity $3$, then $(P',Q') \in \Gamma$ where $P' = \{(x,y) \in D_1^2 : (x,y,x) \in P\}$ and $Q' = \{(x,y) \in D_2^2 : (x,y,x) \in Q\}$. Note that $\Pol(P,Q) \subseteq \Pol(P',Q')$, so adding $(P',Q')$ to $\Gamma$ does not make $\cM_{\SDP} \to \Pol({\Gamma})$ harder to verify. Likewise, since repeating variables is allowed in PCSP instances, whether the basic SDP decides $\PCSP({\Gamma})$ is unaffected by these extra relations.

Assume now that the (alternative) basic SDP decides $\PCSP(\Gamma)$. By Lemma~\ref{lem:free-structure}, it suffices to prove that $\F_{\cM_{\SDP}}(\Gamma)$ is a promise structure. In other words, we must construct a map $\hat{h} : \cM_{\SDP}^{D_1} \to D_2$ such that $\hat{h}(\F_{\cM_{\SDP}}(P_i)) \subseteq Q_i$ for all $i \in [l]$. }

{
By the De Bruijn-Erd{\H o}s Theorem\footnote{{See \cite{ciardo2022clap} or Remark 7.13 of \cite{BBKO21} for additional context.}}~\cite{bruijn1951colour}, it suffices to prove that any finite subsets $F_i \subseteq \F_{\cM_{\SDP}}(P_i)$ for $i \in [l]$, we have that there exists $\hat{h} : \cM_{\SDP}^{D_1} \to D_2$ such that $\hat{h}(F_i) \subseteq Q_i$ for all $i \in [l]$.

Note that each $F_i \subseteq (\cM_{\SDP}^{D_1})^{k_i}$, let $V$ be the subset of $\cM_{\SDP}^{D_1}$ which appears in some tuple of $F_i$ for some $i \in [l].$ It suffices to construct $\hat{h} : V \to D_2$ rather than all of $\cM_{\SDP}^{D_1}$. For each $t \in F := F_1 \cup \cdots F_l$, let $(P_t, Q_t)$ be promise pair corresponding to $t$, and let $S_t \subseteq V$ be the set of distinct elements which appear in $t$. Since $\Gamma$ is closed under identification of variables, we can impose on $t$ a clause $C_t$ corresponding to $(P'_t, Q'_t)$ such that each variable of $S_t$ appears exactly once.

We view $\Phi := (V, \{C_t : t \in F\})$ as an instance of $\PCSP(\Gamma)$. If we can show this instance is accepted by the basic SDP, then we know there is an assignment $\hat{h} : V \to D_2$ which weakly satisfies every clause of $F_1 \cup \cdots \cup F_l$ since replacing $(P_t, Q_t)$ with $(P'_t, Q'_t)$ does not change whether the clause $\hat{h}|_{S_t}(t) \in Q_t$ if and only if $\hat{h}|_{S_t}(S_t) \in Q'_t$. Thus, $\hat{h}$ has precisely the property we want.

Thus, we conclude the proof by constructing an alternative basic SDP solution for $\Phi$. For convenience, let $V = \{x_1, \hdots, x_n\}$ and for each $t \in F$, let $\cF_t$ be the set of possible satisfying assignments $f : S_t \to D_1$ to $C_t$.

For each $i \in [n]$, we let $\bv_i \in \cM_{\SDP}^{D_1}$ be $x_i \in \cM_{\SDP}^{D_1}$ itself. This ensures that (b1) is satisfied. For each $t \in F$, let $\cF_t$ be the set of satisfying assignments to $f : S_t \to D_1$ to $P'_t$. By definition of $\F_{\cM_{\SDP}}(P'_t)$, there exists $\bw_t \in \cM_{\SDP}^{\cF_t}$ such that $\pi_i \oslash \bw_t = x_{i} = \bv_{i}$ whenever $x_i \in S_t$, where $\pi_i : \cF_t \to D_1$ is the evaluation map $\pi_i(f) = f(x_i)$. The condition $\bw_t \in \cM_{\SDP}^{\cF_j}$ is precisely (b2) and the equation $\pi_i \oslash \bw_t = \bv_i$ is precisely (b3) (see Remark~\ref{rem:b3}). Thus, $\Phi$ is indeed accepted by the alternative basic SDP, completing the proof.}
\end{proof}

This completes the proof of Theorem~\ref{thm:basic-SDP}, giving a minion homomorphism characterization of when the basic SDP works for deciding a PCSP.
